\chapter{Results and Discussion}

This chapter presents the experimental findings and performance evaluations of the models trained on the Credit Card Fraud dataset. It details the features implemented in FraudNet, compares the performance of baseline classifiers, gradient boosting models, and the hybrid ensemble, and discusses the mitigation of data leakage and the integration of model interpretability.

\section{Implemented Features}

The FraudNet platform has been successfully realized with several core software and machine learning features:
\begin{itemize}
    \item \textbf{Data Preprocessing \& Cleaning Pipeline}: Automated deduplication, null-value audits, and Robust scaling.
    \item \textbf{High-Performance Serving Endpoints}: A FastAPI backend that accepts Pydantic-validated JSON payloads and returns fraud verdicts along with prediction metadata in less than 50 milliseconds.
    \item \textbf{Interactive Streamlit Web Interface}: A multi-page visualization dashboard featuring:
    \begin{itemize}
        \item \textit{Overview}: Live KPI metrics showing overall accuracy and fraud proportions.
        \item \textit{Score Transaction}: Interactive prediction form that draws samples from the test set or accepts custom inputs, outputting a 5-agent classification trace.
        \item \textit{Model Comparison}: Dynamic ROC/PR curves and comparative metrics.
        \item \textit{Feature Importance}: Global feature impact charts using Random Forest Gini importance.
        \item \textit{Explainability Window}: Local force plots rendering live TreeSHAP values for scored transactions.
    \end{itemize}
\end{itemize}

\section{Results and Performance Analysis}

The model evaluation is conducted on a dataset containing 568,630 transactions. Performance metrics for each classifier under the group-aware train/test partition (80\% training, 20\% test) are reported in Table \ref{tab:model-results-comparison}.

\begin{table}[h!]
    \centering
    \caption{Performance Comparison of Classifiers (Group-Aware Split)}
    \label{tab:model-results-comparison}
    \begin{tabularx}{\textwidth}{| l | c | c | c | c | c |}
        \hline
        \textbf{Model Architecture} 
        & \textbf{Accuracy} 
        & \textbf{Precision} 
        & \textbf{Recall} 
        & \textbf{F1-Score} 
        & \textbf{AUC-ROC} \\ \hline

        Decision Tree (Depth=4) 
        & 0.9493 & 0.9914 & 0.9064 & 0.9470 & 0.9860 \\ \hline

        Logistic Regression 
        & 0.9645 & 0.9770 & 0.9514 & 0.9640 & 0.9933 \\ \hline

        LightGBM (Default) 
        & 0.9993 & 0.9986 & 0.9999 & 0.9993 & 0.9999 \\ \hline

        XGBoost (Default) 
        & 0.9998 & 0.9995 & 1.0000 & 0.9998 & 0.9999 \\ \hline

        Random Forest (n=100) 
        & 0.9999 & 0.9998 & 1.0000 & 0.9999 & 0.9999 \\ \hline

        \textbf{Hybrid RF+RNN Ensemble} 
        & \textbf{0.9979} & \textbf{0.9980} & \textbf{0.9978} & \textbf{0.9979} & \textbf{0.9999} \\ \hline
    \end{tabularx}
\end{table}

The Random Forest and XGBoost classifiers achieved high accuracies and F1-scores, approaching 99.9\% on the tabular split. The custom **Hybrid RF+RNN Ensemble** achieved an F1-score of 0.9979, blending Random Forest ($w_{\gls{rf}} = 0.514$) and SimpleRNN ($w_{\gls{rnn}} = 0.486$) component probabilities. The hybrid ensemble's decision boundary is gated by an optimized classification threshold of $\theta = 0.48$, which balances precision and recall to minimize undetected fraud events.

\section{Comparison with Objectives or Existing Systems}

All initial project objectives were successfully achieved. The system implements a robust training and grid-search pipeline, handles target imbalances using training-fold SMOTE, and deploys both the FastAPI serving layer and Streamlit web app.

Unlike traditional baseline systems that process transactions in isolation, FraudNet captures sequence anomalies. The SimpleRNN sequence component processes transactional history windows of length 6. If a fraud network executes rapid sequential micro-transactions to verify a card's validity, sequence-agnostic classifiers like XGBoost may score each transaction individually as legit because the amount is low. In contrast, the recurrent component in FraudNet detects the high-frequency temporal pattern, resulting in a higher overall fraud risk score.

\section{Challenges and Limitations}

Developing FraudNet presented several key challenges:
\begin{itemize}
    \item \textbf{Information Leakage Prevention}: Applying SMOTE globally across the dataset prior to splitting leads to overfitting and overly optimistic metrics. FraudNet mitigates this by wrapping SMOTE oversampling strictly within training splits, ensuring the test set remains representative of the real-world class distribution.
    \item \textbf{Sequence Construction Overhead}: Formatting structured tabular data into sliding historical windows of shape $(N, 6, 36)$ is computationally expensive. The sequence module optimizes this using contiguous NumPy array striding.
    \item \textbf{SHAP Explainer Compute Time}: Evaluating Shapley values online for deep architectures introduces high inference latencies. FraudNet resolves this by applying TreeSHAP exclusively to the tree-based Random Forest component, keeping the scoring latency under 50ms.
\end{itemize}

\section{Discussion}

The results demonstrate the benefit of combining sequence networks with ensemble models. While classical tree models excel at detecting static outliers, they are blind to sequential dynamics. Conversely, deep recurrent neural networks excel at temporal patterns but require large training budgets. The hybrid ensemble model leverages the strengths of both architectures, achieving an F1-score of 0.9979. Integrating SHAP explanations resolves the black-box nature of the hybrid model, providing risk analysts with clear, real-time feature contributions to verify and audit fraud decisions.
